\documentclass[UTF8,fontset=fandol,zihao=5]{ctexart}

\usepackage[a4paper,margin=2cm]{geometry}
\usepackage{amsmath,amssymb,mathtools}
\usepackage{booktabs,tabularx,array,longtable,multirow}
\usepackage{xcolor}
\usepackage[most]{tcolorbox}
\usepackage{listings}
\usepackage{tikz}

\usetikzlibrary{
  positioning,
  calc,
  decorations.pathreplacing,
  arrows.meta
}

\usepackage{enumitem}
\usepackage{indentfirst}
\usepackage{hyperref}
\usepackage{upquote}

% relationbox 中知识关系图的统一样式
\tikzset{
  relnode/.style={
    font=\small,
    align=center
  },
  relline/.style={
    thick
  },
  relarrow/.style={
    -{Stealth},
    thick
  }
}

\DeclareTextFontCommand{\textbf}{\heiti\color{main}}

\definecolor{main}{HTML}{B85C24}
\definecolor{mainlight}{HTML}{FDF5EF}
\definecolor{method}{HTML}{B88A6E}
\definecolor{methodlight}{HTML}{F8F2EE}
\definecolor{codeback}{HTML}{FBF8F5}
\definecolor{codecomment}{HTML}{7A746F}
\definecolor{codekeyword}{HTML}{9B3F16}
\definecolor{codestring}{HTML}{8A6B32}

\hypersetup{colorlinks=true,linkcolor=main,urlcolor=main}
\setlength{\parindent}{2em}
\linespread{1.5}
\setlength{\abovedisplayskip}{5pt}
\setlength{\belowdisplayskip}{5pt}
\setlength{\abovedisplayshortskip}{3pt}
\setlength{\belowdisplayshortskip}{3pt}
\setlist[itemize]{leftmargin=2em,itemsep=0.1em,topsep=0.2em}
\setlist[enumerate]{leftmargin=2.3em,itemsep=0.1em,topsep=0.2em}

\newfontfamily\headingfont{TeX Gyre Heros}

\ctexset{
  section={format=\centering\heiti\headingfont\zihao{3}\color{main},
           name={第,章},number=\chinese{section},aftername=\quad,
           beforeskip=0.9em,afterskip=0.5em},
  subsection={format=\heiti\headingfont\zihao{5}\color{main},
              name={,、},number=\chinese{subsection},aftername=,
              beforeskip=0.6em,afterskip=0.22em},
  subsubsection={format=\heiti\headingfont\zihao{5}\color{main},
                 number={（\chinese{subsubsection}）},aftername=,
                 beforeskip=0.55em,afterskip=0.15em}
}

\newtcolorbox{relationbox}{
  enhanced,breakable,colback=mainlight,colframe=main,
  boxrule=0.65pt,arc=1.4mm,left=2mm,right=2mm,top=1.2mm,bottom=1.2mm
}
\newtcolorbox{methodbox}[1][]{
  enhanced,breakable,colback=methodlight,colframe=method,
  boxrule=0.65pt,arc=1.4mm,left=2mm,right=2mm,top=1.2mm,bottom=1.2mm,
  title=#1,fonttitle=\kaishu,fontupper=\kaishu,
  before upper={\setlength{\parindent}{2em}}
}

\newtcolorbox{method}[1][]{
  enhanced,breakable,colback=methodlight,colframe=method,
  boxrule=0.65pt,arc=1.4mm,left=2mm,right=2mm,top=1.2mm,bottom=1.2mm,
  title=#1,fonttitle=\bfseries
}

\lstdefinestyle{cppstyle}{
  language=C++,
  basicstyle=\ttfamily\small,
  keywordstyle=\color{codekeyword}\bfseries,
  commentstyle=\color{codecomment}\itshape,
  stringstyle=\color{codestring},
  showstringspaces=false,
  keepspaces=true,
  columns=fullflexible,
  breaklines=true,
  breakatwhitespace=false,
  tabsize=2,
  numbers=left,
  numberstyle=\tiny\color{codecomment},
  numbersep=8pt,
  xleftmargin=1.5em
}

\newtcblisting{cppcode}[1][]{
  enhanced,breakable,listing only,listing engine=listings,
  listing options={style=cppstyle},
  colback=codeback,colframe=main!70!black,
  boxrule=0.55pt,arc=1.2mm,left=1mm,right=1mm,top=1mm,bottom=1mm,
  title=#1,fonttitle=\bfseries
}

\newcommand{\code}[1]{\texttt{#1}}
\newcommand{\cpp}{C\texttt{++}}

\begin{document}

\begin{center}
  {\heiti\headingfont\zihao{1}\color{main} C++ 语言程序设计笔记}
\end{center}

\vspace{0.45em}

\begin{relationbox}
\[
  \boxed{\text{语言与程序}}
  \longrightarrow
  \boxed{\text{类型与数据}}
  \longrightarrow
  \boxed{\text{指针与类}}
  \longrightarrow
  \boxed{\text{控制与函数}}
  \longrightarrow
  \boxed{\text{程序设计}}
\]
\end{relationbox}

\section{基础知识}

\subsection{程序设计语言}

程序设计语言分为机器语言、汇编语言和高级语言。

\begin{enumerate}
  \item 源程序：扩展名为 \code{.cpp}；
  \item 目标程序：扩展名为 \code{.obj}；
  \item 可执行程序：扩展名为 \code{.exe}；
  \item 翻译程序：包括汇编程序、编译程序和解释程序，\cpp 属于编译型语言；
  \item 链接程序：一般由 \code{build} 命令调用。
\end{enumerate}

\[
\boxed{\text{编辑源程序}}
\to \text{源程序}
\to \boxed{\text{翻译}}
\to \text{目标程序}
\to \boxed{\text{链接}}
\to \text{可执行程序}
\]

\subsection{程序组成}

\begin{enumerate}
  \item 注释：\code{//} 表示单行注释，\code{/* */} 表示跨行注释；
  \item 编译预处理：以 \code{\#} 开头，\code{include} 后接头文件；
  \item 程序主体：\code{using namespace std;} 属于程序主体。
\end{enumerate}

\begin{cppcode}[程序基本结构]
#include <iostream>

using namespace std;

int main() {
    return 0;
}
\end{cppcode}

编译顺序：预编译器处理所有编译预处理指令 $\to$ 找到 \code{main()} 函数 $\to$ 编译。

\subsection{内存}

\[
1\ \text{Byte}=8\ \text{bit},\qquad
1\ \text{KB}=1024\ \text{B},\qquad
1\ \text{MB}=1024\ \text{KB}.
\]

\begin{center}
\begin{tabularx}{0.94\textwidth}{>{\bfseries}p{2.8cm}X}
\toprule
分区 & 内容与特点 \\
\midrule
代码区 & 程序指令 \\
栈区（堆栈区） & 局部变量、自动变量、局部常量；先入后出 \\
全局变量区 & 全局变量（\code{extern}） \\
静态变量区 & 静态变量（\code{static}）、静态非基本常量 \\
堆区 & 程序用 \code{new} 申请的空间，申请后必须用 \code{delete} 释放 \\
常量区 & 字符串常量，以及全局常量、静态基本常量 \\
\bottomrule
\end{tabularx}
\end{center}

寄存器型变量 \code{register} 存放于 CPU 寄存器；数量过多时，编译器会自动改按自动变量处理。全局变量与静态变量具有静态生存期并初始化为 $0$；局部变量具有动态生存期，建议显式初始化。

\subsection{程序设计思想}

\begin{center}
\begin{tabularx}{0.94\textwidth}{>{\bfseries}p{2.7cm}X X}
\toprule
维度 & 面向过程 & 面向对象 \\
\midrule
基本结构 & 顺序、选择、循环 & 类与对象；数据成员与成员函数 \\
特性 & 模块化；自顶而下；逐步求精 & 封装、继承、多态 \\
不足或意义 & 数据与过程分离，代码重用不足 & 模块化、软件复用 \\
主要步骤 & 程序分解 & OOA $\to$ OOD $\to$ OOP \\
\bottomrule
\end{tabularx}
\end{center}

程序设计应关注可靠性、成本效益、可理解性和可维护性。

\subsection{词法标记}

\begin{enumerate}
  \item 字符集：大小写英文字母、数字字符和特殊字符；
  \item 关键字：\cpp 保留字；
  \item 标识符：不能是关键字；首字符只能是字母或下划线；其余字符可为字母、下划线或数字；大小写敏感；不得含空格或特殊字符；
  \item 分隔符：\code{( ) [ ] \{ \} < > , : ; ' " // /* */}；
  \item 空白：制表符、空格和空行。
\end{enumerate}

\subsection{运算符与表达式}

\begin{method}[运算符优先级]
从高到低可记为：一元 $\to$ 算术 $\to$ 移位 $\to$ 关系 $\to$ 位 $\to$ 逻辑 $\to$ 三元 $\to$ 赋值 $\to$ 逗号。

一元、三元和赋值运算符通常从右向左结合，其他运算符通常从左向右结合。按“最大运算符”从左向右划分，例如 \code{a+++b} 等价于 \code{(a++) + b}。
\end{method}

\begin{itemize}
  \item 赋值：区分 \code{=} 与 \code{==}；赋值运算符包括一般赋值和复合赋值；
  \item 自增：区分 \code{i++} 与 \code{++i}；
  \item 除法：两侧均为整型时作整数除法，否则作一般除法；注意验零；
  \item 取余：\code{\%} 只能作用于整型，余数符号由被除数决定；
  \item 关系：优先级为 \code{> < >= <=} 高于 \code{== !=}；
  \item 逻辑：\code{!} 高于 \code{\&\&}，\code{\&\&} 高于 \code{||}；存在短路求值；
  \item 条件运算符：\code{?:}；
  \item 逗号运算符：无赋值或自增等副作用的子表达式不改变变量值；
  \item 溢出：按位数截断，可理解为对 $2^{\text{位数}}$ 取余。
\end{itemize}

\begin{cppcode}[整数除法与取余]
int a = -4, b = 3;
std::cout << a / b << std::endl;          // -1
std::cout << 2 * b / a << std::endl;      // -1
std::cout << (float)a / b << std::endl;   // -1.333333
std::cout << a % b << std::endl;          // -1
std::cout << b % a << std::endl;          // 3
\end{cppcode}

表达式由运算符、操作数和分隔符组成；语句通常为表达式加分号；块是由大括号括起的多条语句。

\subsection{类型转换}

\begin{itemize}
  \item 隐式转换：二元运算符两侧类型按类型提升统一；逻辑型中 $0$ 为 \code{false}，非 $0$ 为 \code{true}；赋值时右值转为左值类型；
  \item C 风格：\code{(类型名)表达式}；
  \item \cpp 风格：\code{static\_cast<类型名>(表达式)}。
\end{itemize}

\begin{cppcode}[强制类型转换]
int i = 97;
std::cout << (char)i << std::endl;
std::cout << static_cast<char>(i) << std::endl;

float f = 3.15;
std::cout << (int)f << std::endl;
std::cout << static_cast<int>(f) << std::endl;
\end{cppcode}

\subsection{编译预处理}

\begin{cppcode}[头文件与宏]
#include <iostream>       // 系统头文件
#include "my_header.h"    // 用户头文件

#define M 25
#define MIN(a, b) ((a) < (b) ? (a) : (b))
#undef M
\end{cppcode}

\begin{cppcode}[条件编译]
#ifdef DEBUG
// 调试信息
#endif

#ifndef MY_HEADER_H
#define MY_HEADER_H
// 头文件内容
#endif
\end{cppcode}

\section{变量与常量}

\subsection{基本数据类型}

基本类型包括整型 \code{int}、字符型 \code{char}、实型 \code{float}/\code{double} 和逻辑型 \code{bool}。

\begin{center}
\begin{tabularx}{0.88\textwidth}{>{\bfseries}p{3cm}cX}
\toprule
类型 & 常见长度（字节） & 常见取值范围 \\
\midrule
\code{char} & 1 & $-2^7\sim 2^7-1$ \\
\code{short int} & 2 & $-2^{15}\sim 2^{15}-1$ \\
\code{(long) int} & 4 & $-2^{31}\sim 2^{31}-1$ \\
\code{float} & 4 & 约 $-3.4\times10^{38}\sim 3.4\times10^{38}$ \\
\code{(long) double} & 8 & 约 $-1.7\times10^{308}\sim 1.7\times10^{308}$ \\
\bottomrule
\end{tabularx}
\end{center}

\begin{itemize}
  \item \code{short} 只修饰 \code{int}；\code{long} 可修饰 \code{int} 和 \code{double}；
  \item \code{unsigned} 的取值从 $0$ 开始；\code{long long int} 常占 $8$ 字节；
  \item 实际长度与范围以编译环境为准，可用 \code{sizeof} 检查。
\end{itemize}

\subsection{声明与初始化}

\begin{cppcode}[变量]
int value;
int a = 1;
int b(2);

typedef unsigned short int UINT16;
UINT16 index = 0;

std::cout << sizeof(int) << std::endl;
\end{cppcode}

\subsection{常量}

符号常量必须在声明时初始化：

\begin{cppcode}
const int MAX_SIZE = 100;
\end{cppcode}

\begin{enumerate}
  \item 整型常数：十进制；二进制；八进制以 \code{0} 开头；十六进制以 \code{0x} 或 \code{0X} 开头；后缀可用 \code{L}、\code{U}；
  \item 实型常数：小数形式或指数形式；默认 \code{double}，后缀可用 \code{L}、\code{F}；
  \item 字符常数：\code{'\textbackslash a'} 响铃、\code{'\textbackslash n'} 换行、\code{'\textbackslash t'} 制表、\code{'\textbackslash ooo'} 八进制、\code{'\textbackslash xhh'} 十六进制。
\end{enumerate}

\section{自定义数据类型}

\subsection{字符串}

\subsubsection{C 风格字符串}

\begin{center}
\begin{tabularx}{0.94\textwidth}{>{\bfseries}p{2.2cm}X X}
\toprule
操作 & 字符数组 & 字符指针 \\
\midrule
初始化 & \code{char sArray[] = "This is a book";} & \code{char* sPt = "This is a book";} \\
直接输出 & \code{cout << sArray} & \code{cout << sPt} \\
直接输入 & \code{cin >> sArray} & 不可直接输入 \\
直接修改 & 不可整体改写 & \code{sPt = "OK"} \\
赋值 & 不可整体赋值 & \code{sPt = sArray} \\
\bottomrule
\end{tabularx}
\end{center}

\begin{center}
\begin{tabularx}{0.96\textwidth}{>{\bfseries}p{2.2cm}p{5.6cm}p{2.2cm}X}
\toprule
功能 & 函数原型 & 返回值 & 注意 \\
\midrule
长度 & \code{int strlen(const char* s)} & 长度 & \code{'\textbackslash0'} 不计入 \\
复制 & \code{char* strcpy(char* s1, const char* s2)} & \code{s1} & \code{s1} 空间足够 \\
定长复制 & \code{char* strncpy(char* s1, const char* s2, int n)} & \code{s1} & \code{s1} 空间足够 \\
比较 & \code{int strcmp(const char* s1, const char* s2)} & 字符差值 & 从开头比较 \\
连接 & \code{char* strcat(char* s1, const char* s2)} & \code{s1} & \code{s1} 空间足够 \\
\bottomrule
\end{tabularx}
\end{center}

\subsubsection{C++ 风格字符串}

\begin{cppcode}[std::string]
std::string str1;
std::string str2("Hello!");
std::string str3 = str2;

std::cin >> str1;
std::cout << str1;
std::getline(std::cin, str1);

str1 = "Hello!";
char ch1 = str1[1];
char ch2 = str1.at(1);
std::size_t n = str1.size();
\end{cppcode}

字符串用关系运算符按字符的 ASCII 码依次比较；前部完全相同时，较长者大。

\begin{cppcode}[字符串与数字]
std::string s = "123";
int number = std::stoi(s);
std::string text = std::to_string(123);
\end{cppcode}

\subsection{数组}

声明形式：\code{类型标识符 数组名[常量表达式]}。数组名不可被赋值，访问时必须检查越界。

\begin{cppcode}[数组初始化]
const int NUM = 5;
int scores[NUM] = {70, 100};
int zeros[NUM] = {0};
\end{cppcode}

\begin{cppcode}[字符数组]
const int LEN_STR = 13;
char array1[LEN_STR] = {"Hello World!"};
char array2[LEN_STR] = "Hello World!";
char array3[LEN_STR] = {
    'H', 'e', 'l', 'l', 'o', ' ',
    'W', 'o', 'r', 'l', 'd', '!', '\0'
};
\end{cppcode}

字符数组结尾必须有 \code{'\textbackslash0'}，并为其留出空间。多维数组仅第一维可省略大小，其余维度必须给出。

\subsection{枚举}

\begin{cppcode}[枚举]
enum weekday {sun, mon, tues, wed, thurs, fri, sat};
weekday day = sun;

enum weekday2 {sun2 = 7, mon2 = 1, tues2, wed2,
               thurs2, fri2, sat2};
day = (weekday)3;
int number = sun;
\end{cppcode}

枚举元素本身是常量，默认从 $0$ 开始递增；可指定起始值；可进行关系运算；枚举值赋给枚举变量时需要强制类型转换；通常按 \code{int} 存储。

\subsection{结构与联合}

\begin{cppcode}[结构]
struct student {
    long num;
    char name[20];
};

student s1{}, s2{};
s2 = s1;
\end{cppcode}

结构成员用圆点运算符 \code{.} 访问；结构指针用箭头运算符 \code{->} 访问。只有结构类型相同的变量才能直接相互赋值。

\begin{cppcode}[联合]
union uarea {
    char c_data;
    short s_data;
    long l_data;
};
\end{cppcode}

联合各成员共享存储空间，其大小满足最大成员的存储与对齐要求。联合可无名称，作为无名联合使用。

\section{指针与引用}

\subsection{指针}

\begin{cppcode}[声明与初始化]
int *pa1, *pa2;
char *pch1, *pch2;

int *pa = 0;
double *pb = NULL;
char *pc = nullptr;
\end{cppcode}

指针保存地址；指向对象的解释方式由指针类型决定。未确定指向时应初始化为空指针，\cpp 中推荐 \code{nullptr}。

\begin{itemize}
  \item 一元运算：\code{*} 解引用、\code{\&} 取地址、\code{++} 和 \code{--} 移动指针；
  \item \code{++}/\code{--} 优先级高于 \code{*} 和 \code{\&}，故 \code{*p++} 等价于 \code{*(p++)}；
  \item 指针加减按所指类型的大小移动；
  \item 两指针不能相加，但同一数组中的两指针可相减，结果表示相隔的元素个数。
\end{itemize}

\subsection{动态内存}

\begin{center}
\begin{tabularx}{0.94\textwidth}{>{\bfseries}p{2.2cm}X X}
\toprule
语言 & 申请 & 释放 \\
\midrule
C & \code{malloc(sizeof(类型))} & \code{free(ptr)} \\
\cpp & \code{new 类型(初值)}；\code{new 类型[大小]} & \code{delete ptr}；\code{delete[] ptr} \\
\bottomrule
\end{tabularx}
\end{center}

\begin{cppcode}[动态内存]
int* pn = new int(10);
int* parr = new int[10];

delete pn;
delete[] parr;
\end{cppcode}

\subsection{引用}

\begin{cppcode}[引用]
int someInt = 0;
int &refInt = someInt;
\end{cppcode}

引用必须在声明时初始化；引用关系定义后不可改指向其他对象。

\subsection{多级指针}

取地址运算符 \code{\&} 与解引用运算符 \code{*} 可视为互逆操作，例如 \code{\&*pa} 等价于 \code{pa}。

\begin{method}[数组指针分析]
可借用“微分/积分”记法辅助判断层级：
\begin{itemize}
  \item 一般变量视为原函数，一级指针视为一阶微分，二级指针视为二阶微分；
  \item \code{\&} 视为微分算子，\code{*} 视为积分算子；
  \item 定义数组或指针时，\code{[]} 提升一层；使用时，\code{[]} 降低一层；
  \item $n$ 维数组可按 $n$ 级指针理解；指针移动的单位由解引用后的对象类型决定。
\end{itemize}
该记法主要用于数组指针和多级指针读题。
\end{method}

\begin{cppcode}[多级指针例]
int m[2][3] = {7};
int* pm[] = {
    m[1] - 1,
    &m[1][1],
    &m[0][0] + sizeof(m[1]) / sizeof(int),
    *m + 1,
    *(m + 1)
};

for (int i = 0; sizeof(int*) * i < sizeof(pm); i++) {
    *pm[i] += i + 1;
}
\end{cppcode}

\subsection{指针、引用与函数}

指针和引用均可作为函数形参，通过地址传递修改实参。指针是 C 风格写法，引用是 \cpp 风格写法；编写 \cpp 时通常推荐引用。

\subsection{const 与指针}

\begin{center}
\begin{tabularx}{0.94\textwidth}{>{\bfseries}p{2.8cm}p{5.5cm}X}
\toprule
名称 & 声明 & 含义 \\
\midrule
常指针 & \code{const int *p1} & 所指对象不可经该指针修改，指针可改指向 \\
指针常量 & \code{int *const p2} & 指针不可改指向，所指对象可修改 \\
常指针常量 & \code{const int *const p3} & 所指对象与指向均不可修改 \\
常引用 & \code{const int \&ref} & 不可经该引用修改对象 \\
\bottomrule
\end{tabularx}
\end{center}

\subsection{指针与结构体}

\begin{cppcode}[结构体指针]
struct student {
    int number;
};

student stu{};
student* ps = &stu;
student* dynamicStudent = new student;

ps->number = 1;
delete dynamicStudent;
\end{cppcode}

链表的基本结构是在结构体中保存指向同类节点的指针。

\section{类与继承}

\subsection{类}

类是一组对象的抽象模型。对象的状态对应数据成员，行为对应成员函数。类的成员是抽象描述，对象的成员是具体实例。

\begin{cppcode}[类的声明]
class Student {
public:
    void func();
protected:
    int protectedData;
private:
    int privateData;
};

void Student::func() {
}
\end{cppcode}

成员函数在类内给出函数体时为内联函数；也可在类内声明、类外用作用域运算符 \code{::} 定义。

\subsection{构造函数}

\begin{cppcode}[构造函数]
class Clock {
public:
    Clock(int hour);
    Clock(int hour, int minute = 0);
    Clock(int hour, int minute, int second);
};

Clock::Clock(int hour) {
}
\end{cppcode}

构造函数与类同名，无返回值，通常为公有成员。构造函数可按形参不同重载。若程序未定义构造函数，编译器会补充默认构造函数。

\begin{method}[构造函数歧义]
无参构造函数与“所有参数都有默认值”的构造函数不能同时存在。调用无参构造时二者均可匹配，会产生歧义。
\end{method}

\subsection{复制构造函数}

\begin{cppcode}[复制构造函数]
class Clock {
public:
    Clock(Clock& other);
};

Clock::Clock(Clock& other) {
}
\end{cppcode}

未显式定义时，编译器会生成复制构造函数。若对象内部管理堆空间，则应显式定义复制构造函数，避免多个对象错误地共享同一堆地址。复制构造函数一般不重载。

\subsection{析构函数}

\begin{cppcode}[析构函数]
class Clock {
public:
    ~Clock();
};
\end{cppcode}

析构函数以 \code{\~} 开头，无返回值，不能重载。

\subsection{继承}

继承可分为单继承和多继承。基类也称父类或超类，派生类也称子类。

\begin{cppcode}[继承]
class TEllipse : public TShape {
};
\end{cppcode}

派生类继承基类成员，但不继承构造函数、析构函数和默认赋值运算符。

\begin{center}
\begin{tabularx}{0.94\textwidth}{>{\bfseries}p{2.5cm}ccc}
\toprule
访问属性 & 外部对象 & 派生类成员函数 & 本类成员函数 \\
\midrule
\code{public} & 可 & 可 & 可 \\
\code{protected} & 不可 & 可 & 可 \\
\code{private} & 不可 & 不可 & 可 \\
\bottomrule
\end{tabularx}
\end{center}

\begin{center}
\begin{tabularx}{0.94\textwidth}{>{\bfseries}p{2.6cm}ccc}
\toprule
继承方式 & 基类 \code{public} & 基类 \code{protected} & 基类 \code{private} \\
\midrule
\code{public} 继承 & \code{public} & \code{protected} & 不可访问 \\
\code{protected} 继承 & \code{protected} & \code{protected} & 不可访问 \\
\code{private} 继承 & \code{private} & \code{private} & 不可访问 \\
\bottomrule
\end{tabularx}
\end{center}

\begin{method}[区分访问属性与继承方式]
先判断题目给的是成员的访问属性，还是类声明中的继承方式。访问属性决定“谁能访问”，继承方式决定基类成员进入派生类后的属性。
\end{method}

派生类中的同名成员会覆盖基类同名成员。

\subsection{派生类构造与析构}

\begin{enumerate}
  \item 构造顺序：先基类，后派生类；
  \item 若基类没有无参构造函数，派生类必须在初始化列表中显式调用基类构造函数；
  \item 基类构造函数通常为 \code{public}，也可为 \code{protected}；
  \item 析构顺序：先派生类，后基类。
\end{enumerate}

\begin{cppcode}[派生类构造函数]
class TEllipse : public TShape {
public:
    TEllipse(int longR, int shortR, int x, int y);
    ~TEllipse();
};

TEllipse::TEllipse(int longR, int shortR, int x, int y)
    : TShape(x, y) {
}
\end{cppcode}

\subsection{调用场景}

\begin{center}
\begin{tabularx}{0.92\textwidth}{>{\bfseries}p{3cm}X}
\toprule
函数 & 常见调用场景 \\
\midrule
构造函数 & 创建对象；动态分配；隐式转换 \\
复制构造函数 & 初始化新对象；函数按值传递；返回对象 \\
析构函数 & 形参销毁；临时变量销毁；\code{delete}；对象离开作用域 \\
\bottomrule
\end{tabularx}
\end{center}

\begin{cppcode}[构造、复制与析构]
class Simple {
public:
    int ID;

    Simple(int n) {
        std::cout << "Constructing Simple " << n << std::endl;
        ID = n;
    }

    Simple(Simple &s) {
        ID = s.ID;
        std::cout << "Copy Constructing Simple " << s.ID << std::endl;
    }

    ~Simple() {
        std::cout << "Destructing Simple " << ID << std::endl;
    }

    int getID() {
        return ID;
    }
};

void fun1(Simple s) {
    s.ID = 3;
}

Simple fun2(const int n) {
    Simple simple4(n);
    simple4.ID = 5;
    return simple4;
}

int main() {
    Simple simple1(1);
    Simple simple2 = simple1;
    simple2.ID = 2;
    fun1(simple2);
    fun2(4);

    Simple *pSimple = new Simple(6);
    std::cout << pSimple->getID() << std::endl;
    delete pSimple;
    return 0;
}
\end{cppcode}

原笔记对该程序的计数为：构造函数 $3$ 次、复制构造函数 $3$ 次、析构函数 $6$ 次；实际次数可能受编译器的返回值优化影响。

\section{控制语句}

控制结构分为顺序、选择和循环。算法可用伪代码或流程图表示；流程图中圆角矩形表示开始与结束，矩形表示执行，菱形表示判断，圆形表示连接。

\subsection{选择结构}

\begin{cppcode}[if--else]
if (condition1) {
    // block 1
} else if (condition2) {
    // block 2
} else {
    // block 3
}
\end{cppcode}

\code{else} 与最近的、尚未配对的 \code{if} 配对。

\begin{cppcode}[switch]
switch (expression) {
case constant1:
    // block 1
    break;
case constant2:
    // block 2
    break;
default:
    // default block
}
\end{cppcode}

\begin{itemize}
  \item \code{switch} 表达式只能是字符型、整型或枚举型；
  \item \code{case} 后只能接字符型、整型或枚举型常量；
  \item 匹配后持续执行，直到遇到 \code{break} 或语句结束；
  \item 多个 \code{case} 执行同一段代码时，可连续写标签，只在末尾写一次 \code{break}；
  \item 所有 \code{case} 均不匹配时执行 \code{default}，没有 \code{default} 则跳过。
\end{itemize}

\subsection{循环结构}

\begin{cppcode}[三种循环]
while (condition) {
    // loop body
}

do {
    // loop body
} while (condition);

for (initialization; condition; update) {
    // loop body
}
\end{cppcode}

\code{for} 的执行顺序为：初始化 $\to$ 条件 $\to$ 循环体 $\to$ 更新 $\to$ 条件 $\to\cdots$。三个表达式均可省略，但两个分号不可省略。

\begin{itemize}
  \item 省略初始化：需在循环外赋初值；
  \item 省略条件：形成无限循环，除非通过 \code{break} 等跳出；
  \item 省略更新：可在循环体内修改循环变量；
  \item 省略初始化和更新：形式上接近 \code{while}。
\end{itemize}

\code{break} 跳出循环或结束 \code{switch}；\code{continue} 跳过本轮剩余语句，进入下一轮。

\section{函数}

函数由接口与函数体组成。接口包括函数名、返回值类型和形式参数表；形式参数表由参数类型、个数和顺序决定。

\subsection{声明、定义与调用}

\begin{cppcode}[函数]
int add(int a, int b);      // 声明

int add(int a, int b) {     // 定义
    return a + b;
}
\end{cppcode}

函数调用可作为函数语句、函数表达式或另一个函数的实参。实现调用的函数称主调函数，被调用的函数称被调函数。

\subsection{作用域与可见性}

\begin{enumerate}
  \item 全局作用域，也称文件作用域；
  \item 块作用域，从标识符定义处到本块结束；
  \item 函数作用域；
  \item 类作用域。
\end{enumerate}

内层标识符与外层同名时，内层屏蔽外层。可用全局作用域运算符 \code{::} 访问全局变量。

\begin{cppcode}[全局作用域运算符]
#include <iostream>

double pi = 3.14;

int main() {
    double pi = 3.1415926;
    std::cout << pi << std::endl;
    std::cout << ::pi << std::endl;
    return 0;
}
\end{cppcode}

\subsection{内联、递归与重载}

\begin{itemize}
  \item 内联函数：用 \code{inline} 定义，应先定义再使用；原笔记要求其中不含循环或 \code{switch}；
  \item 递归函数：直接或间接调用自身；必须具有终止条件；
  \item 函数重载：同名函数的参数个数、类型或顺序不同；不能只靠返回值类型区分。
\end{itemize}

\section{编程技巧}

\subsection{输入输出格式}

\subsubsection{C 风格}

\begin{cppcode}[C 风格输入输出]
printf("number: %04d, OK", 33);
putchar('c');

scanf("%d %c %f", &a, &b, &c);
char ch = getchar();
\end{cppcode}

\code{printf} 与 \code{scanf} 的首个参数为格式字符串。多个输入占位符可用空格或逗号分隔。

\subsubsection{C++ 风格}

\begin{center}
\begin{tabularx}{0.96\textwidth}{>{\bfseries}p{6.5cm}X X}
\toprule
函数 & 作用 & 注意 \\
\midrule
\code{std::setw(int n)} & 设置域宽 & 只影响后面一项 \\
\code{std::setiosflags(std::ios::left)} & 左对齐 & 持续生效 \\
\code{std::setiosflags(std::ios::right)} & 右对齐 & 持续生效 \\
\code{std::setfill(char ch)} & 设置填充字符 & 持续生效 \\
\code{std::setprecision(int n)} & 设置浮点有效数字 & 与 \code{fixed} 配合时控制小数位数 \\
\code{std::setiosflags(std::ios::fixed)} & 固定小数形式 & 配合 \code{setprecision} \\
\bottomrule
\end{tabularx}
\end{center}

读取带空格的一整行：

\begin{cppcode}
std::getline(std::cin, str);
\end{cppcode}

\subsection{磁盘文件}

文件可作为输入文件或输出文件。键盘是输入，显示器是输出，磁盘文件可兼具二者。文件分为文本文件和二进制文件。

\begin{cppcode}[文件流]
#include <fstream>

std::ifstream input("input.txt");
std::ofstream output("output.txt");
\end{cppcode}

\code{ifstream} 提供读操作，\code{ofstream} 提供写操作。

\subsection{随机数}

\begin{cppcode}[随机数]
#include <cstdlib>
#include <ctime>

int main() {
    std::srand((unsigned)std::time(NULL));
    int randomNumber = std::rand() % 6;
    return 0;
}
\end{cppcode}

对 $n$ 取余可得到 $0\sim n-1$ 范围内的数。

\subsection{排序}

\begin{cppcode}[标准库排序]
#include <algorithm>
#include <functional>
#include <vector>

std::vector<int> arr;
std::sort(arr.begin(), arr.end());
std::sort(arr.begin(), arr.end(), std::greater<int>());
\end{cppcode}

\begin{cppcode}[快速排序]
void quickSort(int begin, int end) {
    if (begin >= end) return;

    int i = begin, j = end;
    int x = arr.at(begin);

    while (i < j) {
        while (i < j && arr.at(j) >= x) j--;
        if (i < j) arr.at(i++) = arr.at(j);

        while (i < j && arr.at(i) < x) i++;
        if (i < j) arr.at(j--) = arr.at(i);
    }

    arr.at(i) = x;
    quickSort(begin, i - 1);
    quickSort(i + 1, end);
}
\end{cppcode}

\begin{cppcode}[冒泡排序]
const int arraySize = 100;

void bubbleSort() {
    int pos = arraySize - 1;
    while (pos) {
        int bound = pos;
        pos = 0;
        for (int i = 0; i < bound; i++) {
            if (a[i] > a[i + 1]) {
                int hold = a[i];
                a[i] = a[i + 1];
                a[i + 1] = hold;
                pos = i;
            }
        }
    }
}
\end{cppcode}

\subsection{检查清单}

\begin{method}[程序检查]
\begin{enumerate}
  \item 检查边界条件：$0$、$1$、最大值、最小值、负数；
  \item 和变量初始化为 $0$；
  \item 除法前检查除数是否为 $0$，需要小数时使用浮点类型；
  \item \code{switch} 中检查 \code{break}；
  \item 浮点比较可先按需要缩放为整数，减少误差影响；
  \item 检查输出格式和末尾是否多空格；
  \item 没有思路时先列数学模型、变量关系和伪代码；
  \item 读取整行（含空格）时使用 \code{getline} 或 \code{cin.getline}；
  \item 大小写字母 ASCII 码通常相差 $32$，也可用 \code{tolower}/\code{toupper}。
\end{enumerate}
\end{method}

\begin{method}[段错误检查]
常见原因：访问空指针、访问已释放内存、数组越界、修改字符串常量。
\begin{enumerate}
  \item 检查所有指针是否初始化；
  \item 检查数组访问是否越界；
  \item 检查是否尝试修改字符串字面量，尤其是字符指针；
  \item 检查递归是否有终止条件；
  \item 检查内存是否重复释放。
\end{enumerate}
\end{method}

\end{document}
